% 2009 年辽宁卷理科第 15 题：某几何体的三视图（长度单位 m）。
% 来源 PDF：普通高考/2009/2009辽宁理.pdf，SHA256 5ae9708acfe7def5bddd60a34d1474cfa2e2c20a39d130255e0e4393093accad
% 原页：第 1 页（A3 三栏）第 3 栏上部，印刷页 537。
% 重绘依据：正视图为底边 4（左 2 右 2）、高 2 的三角形，自底边中点（x=2）向顶点（x=3）连一条棱；
%   侧视图为高 2、底 3 的直角三角形（竖直边在左）；俯视图为底边 4（左 3 右 1）、
%   深 3 的三角形，自底边 x=3 处向下方顶点（x=2）连一条棱。三视图对应同一个三棱锥，
%   底面三角形底 4 高 3、锥高 2。尺寸按原图标注。
\begin{tikzpicture}[scale=0.9, semithick, >=Stealth,
  every node/.style={inner sep=1.5pt, font=\small}]
  % 正视图
  \begin{scope}[shift={(0,0)}]
    \draw (0,0) -- (4,0) -- (3,2) -- cycle;
    \draw (2,0) -- (3,2);
    \draw[<->] (0,-0.55) -- (2,-0.55);
    \draw[<->] (2,-0.55) -- (4,-0.55);
    \node[below] at (1,-0.62) {$2$};
    \node[below] at (3,-0.62) {$2$};
    \node[below] at (2,-1.05) {正（主）视图};
  \end{scope}
  % 侧视图
  \begin{scope}[shift={(6.3,0)}]
    \draw (0,0) -- (3,0) -- (0,2) -- cycle;
    \draw[<->] (-0.55,0) -- (-0.55,2);
    \node[left] at (-0.62,1) {$2$};
    \draw[<->] (0,-0.55) -- (3,-0.55);
    \node[below] at (1.5,-0.62) {$3$};
    \node[below] at (1.5,-1.05) {侧（左）视图};
  \end{scope}
  % 俯视图
  \begin{scope}[shift={(0,-4.6)}]
    \draw (0,0) -- (4,0) -- (2,-3) -- cycle;
    \draw (3,0) -- (2,-3);
    \draw[<->] (0,0.55) -- (3,0.55);
    \draw[<->] (3,0.55) -- (4,0.55);
    \node[above] at (1.5,0.6) {$3$};
    \node[above] at (3.5,0.6) {$1$};
    \node[below] at (2,-3.25) {俯视图};
  \end{scope}
\end{tikzpicture}
